% Direct-authority diplomatic French for printed p.145.
\SourcePage{145}

\LayerTarget{tohoku:section:2.4:intro}
\noindent suites spectrales, nous renvoyons à [6, Chap. XV et XVII], en nous
contentant de préciser notre terminologie, et de dégager les cas généraux les
plus utiles pour la suite où l'on peut écrire des suites spectrales.

\LayerTarget{tohoku:def:filtered-object}
Soit $C$ une catégorie abélienne. Soit $A\in C$, une \textit{filtration}
(décroissante) sur $A$ est une famille $(F^n(A))$ de sous-trucs de $A$
($n\in\mathbf Z$) avec $F^n(A)\subset F^{n'}(A)$ si $n\geq n'$, un
\textit{objet filtré} dans $C$ est un objet de $C$ muni d'une filtration. Si
$A,B$ sont deux objets filtrés de $C$, un morphisme $u:A\longrightarrow B$ est
dit compatible avec la filtration si $u(F^n(A))\subset F^n(B)$ pour tout $n$.
Avec cette notion de morphisme, les objets avec filtration dans $C$ forment un
catégorie additive (mais non un catégorie abélienne, car une bijection de cette
catégorie n'est pas nécessairement un isomorphisme !).

\LayerTarget{tohoku:def:associated-graded}
On appelle \textit{gradué associé} à un objet filtré $A$ la famille des
$G^n(A)=F^n(A)/F^{n+1}(A)$, on la dénote par $G(A)$\footnote{Ces définitions ne
sont pas autoduales. On rétablit la dualité en associant à toute ``filtration
décroissante'' de $A$ par des $F^n(A)$, la ``cofiltration décroissante
associée'' par les $F'_n(A):=A/F_{1-n}(A)$. Alors, passant d'une catégorie à la
catégorie duale, ces deux filtrations deviennent respectivement une
cofiltration décroissante et une filtration décroissante associées. Il est
encore commode de poser $F_n(A)=F^{1-n}(A)$,
$F'^n(A)=F'_{1-n}(A)$ (filtration et cofiltration \textit{croissantes}
associées aux filtrations précédentes), et suivant les cas, c'est l'une ou
l'autre parmi quatre filtrations ainsi associées dont la considération est la
plus commode. Posant
$G^n(A)=\mathrm{Coker}(F^n(A)\longrightarrow F^{n-1}(A))$,
$G_n(A)=\mathrm{Ker}(F'_n(A)\longrightarrow F'_{n-1}(A))$, on aura
$G^n(A)=G_{-n}(A)$. Les foncteurs $G^n$ et $G_n$ s'échangent si on passe de
$C$ à la catégorie duale.}. $G(A)$ est un foncteur covariant par rapport à
l'objet filtré $A$.

\LayerTarget{tohoku:def:spectral-sequence}
Une \textit{suite spectrale} dans $C$ est un système
$E=(E_r^{p,q},E^n)$ formé a) d'objets $E_r^{p,q}\in C$ définis pour des
entiers $p,q,r$ avec $r\geq2$ (on pourrait remplacer 2 par un entier $r_0$
quelconque, mais dans les applications seul le cas $r=2$ ou $r=1$ semble
intéressant) ; b) de morphismes
\[
d_r^{p,q}:E_r^{p,q}\longrightarrow E_r^{p+r,q-r+1}
\]
tels que $d_r^{p+r,q-r+1}d_r^{p,q}=0$ ; c) d'isomorphismes
\[
\alpha_r^{p,q}:
\mathrm{Ker}(d_r^{p,q})/\mathrm{Im}(d_r^{p-r,q+r-1})
\longrightarrow E_{r+1}^{p,q};
\]
d) d'objets filtrés $E^n$ dans $C$, définis pour tout entier $n$ ; on suppose
que pour tout couple $(p,q)$ fixé, on ait $d_r^{p,q}=0$ et
$d_r^{p-r,q+r-1}=0$ pour $r$ assez grand, d'où on conclut que $E_r^{p,q}$ est
indépendant de $r$ pour $r$ assez grand, on note cet objet par
$E_\infty^{p,q}$ ; on suppose de plus que pour tout $n$ fixé, $F^p(E^n)$ est
égal à $E^n$ pour $p$ assez petit et égal à 0 pour $p$ assez grand ; enfin on
suppose donnés e) des isomorphismes
\[
\beta^{p,q}:E_\infty^{p,q}\longrightarrow G^p(E^{p+q}).
\]
La famille $(E^n)$ (sans filtrations) est appelée \textit{l'aboutissement de
la suite spectral $E$}.

\LayerTarget{tohoku:def:morphism-spectral-sequence}
Un \textit{morphisme $u$ d'une suite spectrale}
$E=(E_r^{p,q},E^n)$ dans une autre $E'=(E_r'^{p,q},E'^n)$ consiste en un
système de morphismes
$u_r^{p,q}:E_r^{p,q}\longrightarrow E_r'^{p,q}$ et
$u^n:E^n\longrightarrow E'^n$ compatibles avec les filtrations, ces morphismes
étant assujettis à commuter aux morphismes $d_r^{p,q}$, $d_r'^{p,q}$,
$\alpha_r^{p,q}$ et $\beta^{p,q}$. Les suites spectrales dans $C$ forment alors
une catégorie additive (mais bien entendu pas une catégorie abélienne).

\LayerTarget{tohoku:def:spectral-functor}
On appelle \textit{foncteur spectral} un foncteur additif défini sur une
catégorie abélienne, à valeurs dans une catégorie de suites
spectrales\footnote{Il semble que pour tous les foncteurs spectraux connus,
l'aboutissement soit en fait un foncteur cohomologique. Il resterait à examiner
les relations entre les homomorphismes bords dans l'aboutissement, et les
autres constituants du foncteur spectral.}.

\LayerTarget{tohoku:def:cohomological-spectral-sequence}
Une suite spectrale est dite \textit{suite spectrale cohomologique} si
$E_r^{p,q}$ est nul pour $p<0$ et pour $q<0$, alors on a
$E_r^{p,q}=E_\infty^{p,q}$ dès que $r>\mathrm{Sup}(p,q+1)$, $E^n=0$ pour
$n<0$, $F^m(E^n)=0$ si $m>n$, $F^m(E^n)=E^n$ si $m\leq0$.

\TohokuContinuation{146}
